% ============================================================
%  CHAPTER IV
%  The Perspective Basis of Byzantine and Old Russian Painting
%  Source pp. 80–93
% ============================================================
\chapter{The Perspective Basis of Byzantine and Old Russian Painting}
\markboth{Chapter IV. The Perspective Basis of Byzantine and Old Russian Painting}{Chapter IV}

\srcpage{80}

The spatial constructions of medieval painting will be examined through the material of
Old Russian and Byzantine art.  Chapter~IX will show that the results obtained here are
essentially of general applicability.

The geometric peculiarities of the depiction of external space on the picture plane in
medieval art differ markedly from those to which the modern viewer is accustomed, and
are often quite incomprehensible to him.  More careful analysis leads to the conclusion
that these differences arise from several sources, which may provisionally be divided
into two groups: those connected with the perspective system, and those not connected
with it.  In this and the two following chapters, only those peculiarities assigned to the
first group will be considered.

The natural desire to understand the perspective system of Old Russian painting is
impeded by its apparent eclecticism: one can observe here elements of both direct and
reverse perspective, and of axonometry.  Additionally, several~--- sometimes very many~---
positions of the viewpoint and other factors play a role.  When examining the perspective
system, it seems rational first to raise the question of the \emph{perspective basis} of
Old Russian painting.  To reveal this basis, several assumptions must be made.  As
before, we shall assume that the aim is to transmit the visible image of the external
world as faithfully as possible, that the picture is painted from a single viewpoint, and
under the assumption of monocular vision.  The peculiarities of depiction that are
connected with the fact that actual vision is binocular and that in practice a single
viewpoint is by no means obligatory in creating a picture will be examined in the
following chapters.  There too will be considered the influence of the form-constancy
mechanism on the depiction of individual objects, and certain other questions connected
with perspective transformations.  The examination of the problem in the present chapter
will thus proceed under a number of simplifying assumptions that are precisely what
make the search for a perspective basis possible.  It is important to note that the
simplifying assumptions used here coincide with those generally accepted for linear
perspective.  This allows a deeper understanding of precisely what brings the perspective
systems of medieval painting and Renaissance painting together and what separates them.

If one raises the question of what kind of perspective the medieval artist used, logic
suggests it should have been perceptual perspective, since the idea that the external
world should be depicted not as one sees it but as it appears on the retina of the eye
could have arisen only at a sufficiently advanced stage of scientific knowledge.

\srcpage{81}

Byzantine and Old Russian art had as its main goal the anthropomorphic depiction of
the deity and saints, for which the transmission of a very close and shallow ``layer'' of
space was entirely sufficient; in this art there was no interest whatsoever in depicting
distant regions (though antiquity did know landscape painting).  But then the perspective
basis for this type of visual art should be the perceptual perspective of the near
foreground~--- that is, axonometry.

The perspective basis may differ from the full perspective system developed upon it,
and may have been to a greater or lesser extent transformed by inevitable deviations
from the idealized initial scheme.  Nevertheless, the premise that the perspective basis
of medieval painting is axonometry is confirmed by analysis of the relevant pictorial
material.  In this analysis an attempt will be made to trace the influence of this basis
both on the methods of depicting individual objects and on the construction of the
picture as a whole.

The main feature of an axonometric image is the preservation of the property of
parallelism of lines in the image.  Lines that in real space are parallel and that are
situated at small distances from the artist should preserve this property in the picture
too.  This should be particularly evident in those cases where the depicted lines are not
only at a small distance from the artist but also at a small distance from one another.
The primacy of parallelism discussed here is often violated by the artist's use of
so-called reverse perspective (the essence of which will be examined in the following
chapters), yet the axonometric basis of Old Russian and Byzantine painting can be traced
with sufficient confidence.

In the icon ``Nativity of John the Baptist'' (Byzantine, 15th~c.) the images of the
cradle and of the roof of the building standing to the right are exact parallelograms~---
that is, are given in axonometry (Fig.\,27).  In the icon ``Entry into the Temple'' (end
of the 15th~c.) the surface of the footstool is an exact parallelogram; all the edges of
the architectural structure on the left side of the picture are parallel to three main
directions~--- horizontal, vertical, and oblique (slightly more than~45° to the horizon);
the sole exception is the left oblique edge of the upper part of the structure, which
makes a certain angle with the other five analogous edges, producing an effect of
reverse perspective.  The axonometric basis is sometimes indicated by the artist's
known indifference to slight deviations from axonometry~--- now in the direction of
direct, now in the direction of reverse perspective.  This is precisely the character of
the architectural structure on the right (Fig.\,28).  In both figures cited, the axonometric
character of the image is conspicuous.  More often, however, the deviations from strict
axonometricity are more considerable (the causes of these deviations will be discussed
below), and then one can speak only of an axonometric basis of the image that has
undergone strong transformation.

\figblock{27}{``Nativity of John the Baptist.'' Byzantine icon, 15th century.
  State Hermitage Museum.  The cradle and the roof of the building at right are exact
  parallelograms~--- strict axonometric images.}

\figblock{28}{``Entry into the Temple.'' Novgorod icon, end of 15th century.
  State Tretyakov Gallery.  All edges of the architectural structure are parallel to three
  main directions; the indifference to small deviations (sometimes toward direct, sometimes
  toward reverse perspective) is characteristic of an axonometric basis.}

\srcpage{82}

The axonometric principle of parallelism, referred to above, sometimes leads to images
striking from a modern viewpoint.  For instance, if an artist depicts several very closely
situated parallel lines (two lines, say, conveying the thickness of a table-top), then in
the picture they usually preserve the property of parallelism.  This is so natural as to
seem self-evident.  Yet this feature can lead to unexpected results when different objects
are involved.  In the upper part of Fig.\,29 the following is schematically shown: three
rectangular footstools, which are rectangular in reality and situated close together~---
just as they are depicted in the Byzantine miniature of the 11th century
``David~between~Wisdom~and~Prophecy.''  As is clear from the scheme, the artist did
not find it possible to deviate from parallelism in depicting the closely situated parallel
edges, even though they belong to different footstools.  As a result, while the left and
right footstools are given in the reverse perspective familiar in medieval painting, the
middle footstool the artist chose to depict differently, merely so as not to violate the
parallelism of the closely situated edges (the distances between the footstools are much
smaller than their width).  This example shows that the artist, working in the system of
perceptual perspective and using the possibilities of choice of distortions, considered
the parallelism of nearby edges more important than the uniformity of the depictions of
the footstools themselves.

\srcpage{83}

Precisely analogous reasons explain the scheme of depiction of rectangular objects in
the icon ``Seventh Ecumenical Council'' (17th~c.), preserved in the Smolensk Cathedral
of the Novodevichy Convent in Moscow (Fig.\,29, lower part).  Here too the parallelism
of closely situated edges is more important to the artist than the uniformity of the
depictions of the objects themselves.  In contrast to the depictions in the upper part of
the figure, the case under consideration belongs to a comparatively late period.  The
artist attempted to render the rectangular platform and the objects standing on it in
direct perspective (this applies to the platform itself and to the base of the lectern);
however, following tradition, he strove to repeat the parallelism of nearby lines, with
the result that the rectangular tables ended up depicted in reverse perspective.\footnote{%
  The tendency to preserve in the image the parallelism of nearby lines, and the
  disregard for the parallelism of lines sufficiently remote from one another, has a
  psychophysiological basis.  As already noted, the field of vision of the eye that is
  characterized by high resolution is very small.  For this reason, in surveying space
  the eye makes a large number of subconscious movements.  Close parallels are
  perceived with good resolution simultaneously, whereas parallels spaced apart in
  space can be seen with the same resolution only by shifting the line of gaze, and
  sometimes by turning the head.  This leads to a person seeing a sufficiently large
  nearby area of space as a collection of local axonometries.  The properties of the
  picture connected with this circumstance will be discussed in more detail in
  Chapter~VI.}

\figblock{29}{Scheme of the depiction of nearby parallel lines as nearby parallel lines
  in the picture.  \emph{Upper}: three rectangular footstools as in the Byzantine miniature
  ``David between Wisdom and Prophecy'' (11th~c.)~--- parallelism of closely spaced
  edges is preserved even across different objects.  \emph{Lower}: scheme from the icon
  ``Seventh Ecumenical Council'' (17th~c., Smolensk Cathedral, Novodevichy Convent)~---
  the same tendency, even when the artist attempts to use direct perspective for some
  elements.}

\srcpage{84}

The axonometric basis of Old Russian painting is also manifested in the peculiarities of
conveying spatial depth.  In the system of axonometry, when depicting shallow regions
of space, only two monocular depth cues can be used: the first~--- occlusion (nearby
objects block more distant ones)~--- and the second~--- more distant objects appear
closer to the horizon line.  As was already emphasized in Chapter~II, the latter cue is
not absolute, and moreover can almost always be formulated more simply: more distant
objects should be depicted higher than nearer ones.  Possible exceptions, arising for
instance in depicting ceilings, are almost never encountered in medieval art.  Linear
perspective, geometrically speaking, has in addition to the above one further, and for
the system of linear perspective most important, cue~--- the diminution of the visible
sizes of objects as their distance from the picture plane increases.  Since this last cue
is in principle inapplicable in axonometry, the first two become the only, and hence
the principal, ones.  Connected with this is the fact that in Old Russian painting the
vertical displacement of figures to show their distribution in depth in the depicted space
becomes a clear artistic technique.  This displacement is often very considerable and
is expressed in a markedly elevated horizon and~--- even when depicting shallow
spaces~--- is significantly more pronounced than in the system of linear perspective.

\figblock{30}{Scheme of the depiction of two objects one of which blocks the other.
  \emph{Upper}: two parallelepipeds of different height, at the same distance from the
  picture plane, as a modern draftsman would render them.  \emph{Lower}: the same
  group as a medieval artist might prefer~--- the shorter parallelepiped is deformed to
  avoid occlusion, emphasizing the equal distance of both objects from the viewer.}

The first of the depth cues mentioned above~--- occlusion~--- also acquires special
significance.  Connected with this is, in particular, the tendency observable in many
cases for the artist to depict objects equidistant from him in such a way that they do
not block one another, even if this is unavoidable because they differ in height.  In the
upper part of Fig.\,30 is given the image of a group of two parallelepipeds of different
heights (but situated at the same distance from the picture plane) as a modern
draftsman would render it.  In the lower part of the figure is given another image of
the same group, which a medieval artist would prefer.  As was shown above, the
distortion of the visible geometric appearance of an object (for example, of the
footstools and tables in Fig.\,29) for the sake of a more correct rendering of some
``more important property'' (the parallelism of closely situated edges in Fig.\,29) is
entirely acceptable from the standpoint of the medieval master.  Therefore also in the
case depicted in Fig.\,30, distorting the form of the shorter parallelepiped in order to
emphasize the equal distance of both objects from the viewer may prove preferable.

\srcpage{85}

To illustrate what has been said, let us turn to the icon ``Entry into the Temple'' from
the iconostasis of the Cathedral of the Dormition in Kirillov.  The shape of the
footstool on which Mary stands (the strong narrowing of its front part) was dictated
by the desire not to block it with Anna's garment~--- that is, the non-intersection of
the images appeared to the author of the icon more important than the correct rendering
of the footstool's form (Fig.\,31).

\figblock{31}{``Entry into the Temple.'' Icon from the Kirillo-Belozersky Monastery.
  15th century.  Fragment.  The footstool on which Mary stands is strongly narrowed at
  the front to avoid being covered by Anna's garment~--- non-intersection was judged
  more important than correct form.}

Occlusion as the primary depth cue becomes so habitual that it persists even in
17th-century works in which linear perspective has already made an appearance.
Figure~32 gives the scheme of the legs of a throne and of a footstool standing before
it, depicted in the icon ``Old Testament Trinity'' (by Yakov Kazanets and Gavrila
Kondratyev) from the Church of the Trinity in Nikitniki, Moscow.  As is clear from
the scheme, the front legs of the throne and footstool begin at the same level~--- at the
lower edge of the icon~--- yet the footstool stands entirely in front of the throne, as is
evident from its occluding the throne's image.  The artists are untroubled by the
circumstance that from a modern point of view the rear legs of the footstool have
``flown up'' and do not rest on the ground.

\srcpage{86}

Here there is a clear manifestation of the preference for occlusion over other depth
cues~--- even over the rule, known to medieval art, by which a more distant object should
be depicted higher than a less distant one.

\figblock{32}{Image of footstool and throne legs (tracing from icon).  ``Old Testament
  Trinity,'' by Yakov Kazanets and Gavrila Kondratyev.  Church of the Trinity in
  Nikitniki, Moscow.  The footstool's front legs are at the same level as the throne's,
  but the footstool occludes the throne~--- demonstrating the primacy of occlusion as a
  depth cue; the footstool's rear legs appear to float.}

The axonometric basis of the image of the near foreground inevitably leads to the depth
of the depicted space being small.  Indeed, as follows from the material presented in
Chapter~II, axonometry correctly conveys volume only in shallow spaces.  Particular
interest attaches to the methods by which the artist was compelled to limit the depth of
space.

The simplest type of construction of an essentially axonometric shallow space is the
depiction of a saint against a monochrome background with a comparatively narrow
foreground strip.  The axonometric character of such a construction is manifested in the
method of depicting the extended objects included in the composition.  This type is
frequently encountered in early Byzantine art, particularly in the Ravenna mosaics (the
Mausoleum of Galla Placidia; the Basilica of Sant'Apollinare~Nuovo).\footnote{%
  See: V.\,N.~Lazarev. \emph{Byzantine Painting}. Moscow, 1971, pp.\,44, 56 (images
  of the lattice of St.\,Lawrence's fire, the base of the Virgin's throne).}
  It should be noted that sometimes the extended objects described are given not in
strict axonometry but in slight reverse perspective.  This does not contradict the
assertion about the axonometric basis of such painting, since slight reverse perspective
arises (as will be shown in the next chapter) as a result of the transformation of a
fundamentally axonometric image.

Eliminating depth of the depicted space by introducing a monochrome background is
not the only technique, especially for icons of the festival tier, hagiographic scene
panels, and the like.  In such cases the irreal conception of space gives way to a
narratively justified one.  In compositions of this type~--- beginning with the
simplest~--- the space is bounded by some flat element.  For instance, in the miniatures
of the Armenian Gospel of Queen Mlke (12th~c., Venice), evangelists are depicted
against a background of nearby curtains, which immediately sharply limits the depth of
the picture's space.  From a formally geometric standpoint this composition is
conditioned by the artist's aspiration to keep the depth of space within the magnitude
permissible for an axonometric image.  Similarly instructive are the numerous icons of
the ``Crucifixion,'' in which the scene of the crucifixion is given against a background
of the Jerusalem wall.  In both cases described, the curtain or wall justifies the absence
of a depiction of the distance, which, as noted above, coordinates poorly with the
axonometric image of the near foreground.

\srcpage{87}

More often, architectural backgrounds serve to limit the depth of space.\footnote{%
  It is well known that an action depicted in front of a building should often be
  understood as taking place inside it.  For what follows this is inessential, since only
  the formal-geometric properties of the image are discussed below.}
This widely employed motif, among other things, introduced variety into the composition
and simultaneously opened wide scope for the artist~--- the limitation of the depth of
the space in which the action takes place did not hinder, if one may put it so, the
``spatial narrative.''  The techniques employed here did not stimulate the size-constancy
mechanism, since they corresponded to the scheme shown in the lower part of Fig.\,20.
Indeed, the depicted space in this case was divided into two axonometric layers~---
a foreground in which the action takes place, and a background (the architecture)~---
with the transition from one layer to the other being not continuous (as in linear
perspective) but step-like, since the linear scale of the depicted objects changed in a
step-wise manner.  Depending on the composition, this transition was sometimes not
masked at all (in such cases the architectural background resembles a stage backdrop at
the rear of a theatre)~--- or was masked in a manner similar to that of the lower scheme
in Fig.\,20.

The first type of image is encountered very frequently.  Practically all icons of the
``Entry into the Temple,'' ``Presentation in the Temple,'' ``Nativity of the
Theotokos,'' and many others are built on this principle.  As examples of such
compositions one may cite the icons mentioned above (Figs.\,27, 28).  In addition to
architectural backgrounds, other images~--- most often stylized icon rocks~--- were also
used to bound the near axonometric space.  Let us recall the composition ``John the
Theologian and Prochorus on the Island of Patmos.''

\srcpage{88}

Another method of depiction is also known, in which the step-like transition from one
axonometric space to the next was masked in one way or another.  In Andrey Rublev's
``Trinity'' (Fig.\,33) this division of the two planes is achieved with exceptional
tact~--- the nearly merged images of the angels' wings created a kind of ``curtain'' and
concealed the transition from the near foreground to the background, which contains the
ideologically and symbolically important images of a building, the Oak of Mamre, and
a mountain.  Rublev's ``Trinity'' immediately after its creation became a model for
Russian icon-painters.  Among the many icons of this type there are those in which
the wings of the angels do not merge and thus do not screen the transition from the
foreground to the background (see, for example, the ``Trinity'' from the village of
Borodava (Fig.\,34)).  Comparison of such icons with others in which the ``curtain'' of
wings is preserved shows how justified in this composition the masking of the step-like
transition from the foreground to the background is~--- making it possible to avoid the
topographically inexplicable juxtaposition of mountain, building, and altar table.

The need to ``justify'' the break in the scale of the image (in particular by masking)
was already discussed in Chapter~III.  The small depth of the picture's space derived
from the peculiarities of perceptual perspective of the near foreground, whose geometric
properties are close to axonometry.  This small spatial depth imposed its own
constraints on the composition: it required that the action (more precisely, the
interaction of the depicted persons) be structured parallel to the base of the canvas~---
i.e.\ from left to right or vice versa.  Indeed, in Russian medieval painting, when a
single subject is depicted, this condition is always observed.

\figblock{33}{Andrey Rublev. \emph{Trinity.} 1420s.  State Tretyakov Gallery.  The
  nearly merged wings of the three angels form a ``curtain'' that screens the step-like
  scale transition from the near foreground (the figures) to the distant background
  (building, Oak of Mamre, mountain).}

\figblock{34}{``Trinity.'' Icon from the village of Borodava, 15th century.
  Andrey Rublev Museum of Ancient Russian Culture and Art.  The wings do not merge
  and do not screen the foreground-to-background transition; comparison with Rublev's
  version makes clear how much is lost without the masking.}

To illustrate the assertion, we may refer to the icon ``Entry into the Temple'' (Fig.\,28).
Here the procession of Joachim, Anna, Mary, and their companions, and the movement
of the priest of the Jerusalem temple coming to meet Mary, are subject to a direction
parallel to the lower edge of the icon.  Nor could it have been otherwise, since the depth
of the space (without counting the second, architectural, plane) is insufficient to
accommodate so many persons, given the condition that the full figure of each main
personage must be shown.

\srcpage{89}

The difficulty that would arise with a different compositional solution is apparent from
the difficulty and incompleteness with which the artist depicted the secondary
participants in the procession.  Their heads are displaced upward in accordance with
the depth cues mentioned above, and their figures are partly blocked by the main
personages.  The second subject depicted on the same icon~--- ``The Apparition of the
Angel to Mary''~--- not connected with the first either in time or place, and placed on
the second plane, is also subject to the direction from left to right.  Entirely analogous
is the composition of many other icons of the ``Entry into the Temple,'' as well as of
such widely-spread subjects as the ``Presentation in the Temple,'' the ``Raising of
Lazarus,'' the ``Crucifixion,'' and the like.

Particular interest belongs to the study of compositions of icons in which by the logic
of the depicted event it is necessary to show not only the direction parallel to the lower
edge of the icon but also the direction perpendicular to it.

In the icons of the ``Baptism,'' if the action is developed in the direction parallel to the
base of the canvas (as was always done), then the Jordan River must flow toward or
away from the viewer, since traditionally John the Baptist and the angels are placed on
different banks.  The way out of this situation was found by the artist as follows: he
depicted a section of the river situated in the shallow space of the picture, after which
the image of the river was cut off.  To the right and left of this conventional water
horizon, mountains were depicted, and in this way the shallow axonometric space of
the action was clearly delineated.  Sometimes this aspiration to a material bounding of
the space went so far that the iconographer turned the Jordan River into a small pond
bounded by land at what he considered the limit of the space to be depicted on the icon.
This is precisely what the iconographer did in the panel ``Baptism'' of the icon
``Theotokos of Vladimir'' (1514), in the Annunciation Cathedral of the Moscow Kremlin
(Fig.\,35).

\figblock{35}{``Baptism.'' Panel from the icon ``Theotokos of Vladimir,'' 1514.
  Annunciation Cathedral, Moscow Kremlin.  The Jordan River is reduced to a small pond
  bounded by land~--- a solution that keeps the depicted water within the shallow
  axonometric space, eliminating the compositional problem of a river flowing
  perpendicular to the picture plane.}

\srcpage{90}

Such a depiction, from the standpoint of the medieval artist, cannot be denied a certain
logic, as is attested by a literally identical technique in the panel ``Baptism'' of the
icon ``John the Forerunner, Angel of the Desert'' from the Church of Nikola Nadein
(Yaroslavl, 17th~c.); an analogous composition may be seen in the Armenian Gospel
of~1038.\footnote{%
  L.\,A.~Durnovo. \emph{Armenian Miniature}. Yerevan, 1969, Table~6.}
If it seemed important to the iconographer to show the river in a more natural manner,
there were various possible ways.  In Old Russian painting, the aspiration to show the
Jordan in full led to dividing it into two branches that flowed over mountains located
at depth in the axonometric space~--- i.e.\ without increasing the depth of the near
foreground.  This is done, for example, in the Pskov icon ``Baptism.''\footnote{%
  A.\,N.~Ovchinnikov, N.\,V.~Kishilov.  \emph{Painting of Ancient Pskov,
  13th--14th Centuries}. Moscow, 1971, Table~8.}

\srcpage{91}

What has been said testifies to the known compositional difficulties that the Old Russian
artist had to overcome owing to the fact that the method of depiction he employed had
an axonometric basis.

The examples examined above show that the peculiarities of Old Russian painting are in
agreement with the formal geometric requirements dictated by the logic of image
construction in the case of the artist's using perceptual perspective of the near
foreground.  The methods of depicting interiors encountered in Old Russian art say the
same.  Of all the possible schemes for depicting an enclosed space, artists almost always
chose the simplest (Fig.\,18-E)~--- one that allowed them to make do with the far wall
(or some other limitation of the depth of space) and a foreground strip.  In very rare
cases scheme~18-B or compositions inclining toward it are encountered.  Comparatively
late, depictions of interiors built according to scheme~18-C appear~--- for example in
the paintings of the Church of the Trinity in Nikitniki (Moscow).\footnote{%
  E.\,S.~Ovchinnikova. \emph{The Church of the Trinity in Nikitniki}. Moscow, 1970,
  pp.\,55, 63.  It is interesting to note that the aspiration to show the interior fully and
  volumetrically also attracted artists to the variant shown in Fig.\,18-A, which is
  characteristic of the period preceding the Renaissance~--- in particular, it appears
  frequently in Giotto.}
Without examining the succession of techniques in interior depiction, it suffices to
state that the methods of depicting an enclosed space used before artists had mastered
linear perspective correspond to schemes admissible from the standpoint of the theory
of perceptual perspective.

\srcpage{92}

The assertion~--- advanced at the beginning of the present chapter on theoretical
grounds~--- that the perspective basis of Old Russian painting is axonometry is
confirmed by a number of examples.  True, the axonometric character of the perspective
construction in Old Russian art is by no means conspicuous.  This is entirely natural,
since it is a question only of a perspective \emph{basis}, not of a rigorous perspective
system.  Just as the skeleton~--- invisible to the eye~--- is the formative foundation of
the human figure, axonometry is that very same foundation: inconspicuous, yet holding
together and unifying the individual elements of the image into a single whole.  Since
knowledge of the basis does not yet mean knowledge of the subject in its essentials,
the next chapter is devoted to a more detailed analysis of the perspective system in Old
Russian and Byzantine painting.
